\section{Auto-off de overlays}

La funcionalidad \textit{auto-off} resuelve un problema frecuente en producción en
directo: el operador realiza un corte de cámara olvidando retirar el rótulo activo,
que queda superpuesto sobre la nueva fuente de forma inapropiada. La consecuencia
habitual es que el nombre de un ponente aparece sobre la imagen de otro, lo que
resulta en un error visible para el espectador.

\subsection{Funcionamiento}

Cuando la opción \texttt{auto-off = true} está activa en la sección \texttt{[overlay]}
del fichero de configuración, el módulo \texttt{voctocore/lib/overlay.py} registra un
\textit{listener} sobre el evento \texttt{video\_status}. Cada vez que la fuente de
vídeo activa cambia, el \textit{listener} inicia automáticamente la secuencia de
fundido de salida del overlay PNG activo, retirándolo del aire antes de que el nuevo
plano sea completamente visible para el espectador.

Adicionalmente, si hay textos dinámicos superpuestos sobre el rótulo (nombre del
ponente, subtítulos), éstos se eliminan junto al overlay gráfico, garantizando que
ningún elemento de rotulación quede huérfano en pantalla.

\subsection{Configuración}

El comportamiento se controla mediante los siguientes parámetros en la sección
\texttt{[overlay]} del fichero INI:

\begin{itemize}
  \item \texttt{auto-off = true/false} --- activa o desactiva la funcionalidad.
  \item \texttt{user-auto-off = true/false} --- permite al operador activar/desactivar
        el auto-off desde la GUI sin modificar el fichero de configuración.
  \item \texttt{blend-time = <ms>} --- tiempo de fundido de salida en milisegundos.
\end{itemize}

% --- Figura sugerida ---
% \begin{figure}[H]
%   \centering
%   \includegraphics[width=0.75\textwidth]{figuras/secuencia_autooff.png}
%   \caption{Diagrama de secuencia: corte de cámara con auto-off activo.}
%   \label{fig:secuencia_autooff}
% \end{figure}
